\usepackage[english]{babel}
\usepackage{tikz}
\usetikzlibrary{arrows.meta, positioning}
\usepackage{amsmath}
\usepackage{amssymb}
\usepackage{amsfonts}
\usepackage[thmmarks, amsmath]{ntheorem}
\usepackage{setspace}
\usepackage{csquotes}
\usepackage{enumitem}
\usepackage{multicol}
\usepackage{parskip}
\usepackage{microtype}
\usepackage{hyphenat}
\usepackage{booktabs}
\usepackage{array}
\usepackage[margin=1in]{geometry}
\usepackage{fancyhdr}
\usepackage{titlesec}
\usepackage[
    pdfauthor={Lance Pollard},
    pdftitle={Vibe Theory: A Simulable Discrete Universe},
    pdfsubject={Discrete Physics, Causal Sets, Lattice Gauge Theory, Quantum Foundations, Monism},
    pdfkeywords={monism, discrete physics, causal sets, lattice gauge theory, overlap fermions, Bell inequalities, emergent spacetime, vibe theory},
    pdfcreator={Lance Pollard},
    pdfproducer={LuaLaTeX},
    colorlinks=true,
    linkcolor=black,
    urlcolor=black,
    citecolor=black
]{hyperref}

\hypersetup{
    pdfinfo={
        ContactEmail={lancejpollard@gmail.com},
        ContactURL={https://clue.surf},
        Copyright={Copyright (C) 2026 ClueSurf}
    }
}

\usepackage{etoolbox}

\widowpenalty=10000
\clubpenalty=10000

\setlength{\parindent}{0pt}
\setlength{\parskip}{1em}

\pagestyle{fancy}
\fancyhf{}
\renewcommand{\headrulewidth}{0pt}
\fancyfoot[C]{-- \thepage\ --}
\setlength{\footskip}{40pt}

\fancypagestyle{plain}{
  \fancyhf{}
  \fancyfoot[C]{-- \thepage\ --}
  \renewcommand{\headrulewidth}{0pt}
}

\AtBeginEnvironment{equation}{\Large}
\AtBeginEnvironment{align}{\Large}
\AtBeginEnvironment{gather}{\Large}

\let\oldequation\equation
\let\endoldequation\endequation
\renewenvironment{equation}{%
    \noindent\vspace{-\parskip}\vspace{-\baselineskip}%
    \oldequation
}{%
    \endoldequation
    \noindent\vspace{-\parskip}\vspace{-\baselineskip}%
}

\BeforeBeginEnvironment{equation}{\vspace{8pt}}
\AfterEndEnvironment{equation}{\vspace{16pt}}

\setlength{\itemsep}{0pt}
\setlist[itemize]{before=\vspace{0pt},after=\vspace{-2pt}}

\titleformat{\section}
  {\Large\bfseries}
  {\thesection.}
  {0.5em}
  {}

\titleformat{\subsection}
  {\normalsize\bfseries}
  {\thesubsection.}
  {0.5em}
  {}

\titleformat{\subsubsection}
  {\small\itshape}
  {\thesubsubsection.}
  {0.5em}
  {}

\onehalfspacing

\theoremstyle{definition}
\theoremstyle{axiom}
\theoremstyle{theorem}
\theoremstyle{lemma}
\theoremstyle{proposition}
\theoremheaderfont{\bfseries}
\theorembodyfont{\normalfont}
\theoremindent0pt
\newtheorem{definition}{Definition}
\newtheorem{axiom}{Axiom}
\newtheorem{theorem}{Theorem}
\newtheorem{lemma}{Lemma}
\newtheorem{proposition}{Proposition}
\newtheorem{principle}{Principle}
\newtheorem{result}{Result}

\makeatletter
\def\@begintheorem#1#2#3{%
  \par\addvspace{1em}%
  \noindent\textbf{#1 #2.}
  \ifx\@empty#3\relax
    \quad
  \else
    \ \textbf{#3:}
  \fi
  \hspace{0.5em}
}
\def\@endtheorem{\par}
\makeatother

\title{
  {\huge \textbf{Vibe Theory}} \\[2mm]
  \large A Simulable Discrete Universe: \\ How Space, Time, Matter, and Force Can Emerge from One Mesh
}

\author{Lance Pollard \\ \small lancejpollard@gmail.com}
\date{June 6, 2026}
